% !TEX root = ../main.tex
%%
% BIThesis 本科毕业设计论文模板 —— 使用 XeLaTeX 编译
%%

% 第一章节
\chapter{绪论}
\section{研究背景与意义}
软件测试是保证软件质量的关键环节，在软件开发生命周期中占据重要地位。传统软件测试过程高度依赖人工经验，存在效率低、覆盖不足、成本高等问题。据统计，软件测试环节平均消耗整个开发周期40\%-50\%的时间与资源，在安全关键系统中甚至高达60\%-80\%\cite{Myers2011}。

敏捷开发与持续集成模式的应用使代码变更更加频繁，迭代周期显著缩短。回归测试成为保证软件质量的重要手段，但传统回归测试主要依赖人工编写测试用例，耗时耗力且难以快速响应频繁的代码变更。开发人员平均需要花费20\%-30\%的工作时间编写和维护测试用例，且人工测试用例的质量高度依赖测试人员的经验与技能。

代码变更影响分析是回归测试的核心环节，其目标在于识别代码变更可能影响的系统范围，从而指导测试资源的分配与测试用例的生成。传统方法主要基于文本差异比较，仅能识别行级变化，难以反映代码的真实语义与影响范围。例如，当函数内部逻辑发生重构或接口参数发生变化时，仅凭文本差异难以准确识别其对其他模块的影响。

大语言模型（Large Language Model, LLM）与多智能体系统（Multi-Agent System, MAS）的快速发展，为软件测试自动化带来了新机遇。大语言模型在代码理解、知识推理等方面展现出强大能力，多智能体系统通过多个Agent的协同决策与任务分工，能够处理复杂多步骤的任务，提高系统的智能化水平与可扩展性\cite{Wooldridge2002}。

基于上述背景，本课题围绕基于多智能体的软件测试自动化展开研究，旨在解决传统软件测试过程中对人工经验依赖强、测试覆盖不足及效率较低等问题。本研究的意义主要体现在：（1）理论意义：提出软件测试自动化新框架，丰富软件测试智能化研究的理论体系；（2）技术意义：有效提升测试用例生成的自动化程度与智能化水平；（3）应用价值：可应用于持续集成环境下的回归测试场景，支持快速迭代开发模式。

\section{国内外研究现状}
\subsection{软件测试自动化研究现状}
软件测试自动化旨在通过自动化工具与技术，减少或替代人工测试工作，提高测试效率与质量。根据自动化程度的不同，软件测试自动化可分为测试脚本自动化、测试用例生成自动化与测试执行自动化等层次。

在测试用例自动生成方面，传统方法主要分为基于搜索的方法（Search-Based Software Testing, SBST）与基于模型的方法（Model-Based Testing, MBT）\cite{AliBriand2007}。基于搜索的方法通过优化算法在测试用例空间中搜索，以最大化覆盖率或最小化测试数量为目标，代表性工具包括Pynguin、EvoSuite等。Pynguin是针对Python语言的单元测试生成工具，采用遗传算法优化测试用例生成过程，支持多种覆盖准则\cite{LukasczykKF2023}。然而，基于搜索的方法主要关注覆盖率指标，生成的测试用例往往缺乏语义合理性，可读性与可维护性较差。

基于模型的方法通过构建软件系统的形式化模型自动生成测试用例，但需要人工构建模型，工作量大且难以适应频繁的代码变更。

近年来，基于神经网络的测试用例生成方法逐渐受到关注，代表性工作包括A3Test、ATLAS等\cite{Watson2019}。此类方法需要大量标注数据进行训练，且生成的测试用例质量依赖于训练数据的质量，泛化能力有待提升。

大语言模型的出现为测试用例自动生成带来了新可能。GPT-4、Claude等大语言模型在代码理解与生成方面展现出强大能力。研究表明，基于大语言模型的测试生成方法在测试质量与可读性方面优于传统方法，但仍存在生成结果不稳定、缺乏约束与验证等问题\cite{Riedel2023LLMTest}。

\subsection{代码变更影响分析研究}
代码变更影响分析是软件维护与回归测试的重要环节，其目标在于识别代码变更可能影响的系统范围，从而指导测试资源的分配与回归测试的选择\cite{BohnerArnold1996}。

静态分析方法通过分析代码的结构信息推断代码变更的影响范围。代表性工作包括Horwitz等提出的基于系统依赖图（System Dependency Graph, SDG）的影响分析方法\cite{Horwitz1990}和Ren等提出的Chianti\cite{Ren2004}。Horwitz等的方法通过构建系统依赖图，利用控制依赖与数据依赖关系分析代码变更对其他模块的影响传播路径。Chianti通过将代码变更分解为原子变更集合，并结合调用图分析识别受影响的测试用例。静态分析方法的优点在于无需执行程序，分析成本低，但存在误报率较高的问题，难以准确识别动态绑定、反射调用等场景下的影响。

动态分析方法通过执行程序并收集运行时信息识别代码变更的影响范围。代表性工作包括Daikon\cite{Ernst2001Daikon}和Jalangi\cite{Sen2015}。Daikon通过动态检测程序不变量，推断代码变更可能违反的约束条件。Jalangi基于动态分析技术，收集JavaScript程序的执行轨迹，支持变更影响分析与回归测试选择。动态分析方法的优点在于分析结果较为准确，但需要执行测试用例，分析成本较高。

近年来，研究者开始将静态分析与动态分析相结合，构建混合式影响分析方法。Chianti与Ekstazi等工具通过静态分解识别变更影响范围，再通过动态调用图验证影响是否真实发生，有效降低了测试执行成本\cite{Ren2004,Gligoric2015Ekstazi}。然而，现有方法主要关注代码结构层面的影响，对代码语义层面的影响识别能力有限。

\subsection{多智能体与大模型在软件工程中的应用}
多智能体系统由多个智能体组成，各智能体通过协同决策与任务分工，共同完成复杂任务。多智能体系统具有自治性、协作性、适应性等特点，在复杂系统建模、任务规划、决策支持等领域得到广泛应用。

近年来，多智能体系统在软件工程中的应用逐渐受到关注。AutoGen是微软提出的通用多智能体框架，支持多个大语言模型Agent之间的对话与协作，可应用于代码生成、测试、调试等场景\cite{Wu2023AutoGen}。MetaGPT面向软件开发的多智能体系统，通过模拟软件开发团队的角色分工，实现需求分析、设计、编码、测试等任务的自动化\cite{Hong2023MetaGPT}。

大语言模型在软件工程中的应用近年来取得显著进展。在测试生成方面，代表性工作包括CoditT5、ChatTester等\cite{Zhang2022CoditT5}。CoditT5基于CodeT5模型，根据代码变更自动生成测试用例。ChatTester基于ChatGPT，通过多轮对话方式引导测试用例的生成与优化。

然而，现有研究主要将大语言模型作为单一组件使用，缺乏多组件之间的协同机制。如何构建多智能体协同框架，实现各环节之间的信息流转与任务联动，是当前研究的空白。

综上所述，现有研究在测试自动化、影响分析、大语言模型应用等方面取得了一定进展，但仍存在以下不足：（1）传统测试自动化方法主要关注覆盖率指标，缺乏对测试语义合理性的考虑；（2）现有影响分析方法主要关注结构层面，对语义层面的影响识别能力有限；（3）大语言模型在测试生成中的应用主要采用单一模型，缺乏多组件协同机制。针对上述问题，本课题提出基于多智能体的软件测试自动化方法，通过代码语义分析、影响分析、缺陷模式挖掘与大语言模型、多智能体系统的深度融合，实现从代码变更到测试用例生成的端到端自动化。

\section{本文主要研究内容}
本课题围绕基于多智能体的软件测试自动化展开研究，主要研究内容包括以下五个方面：

（1）代码变更分析方法研究。引入抽象语法树（AST）与Tree-sitter技术，对源代码进行语法级解析，将非结构化代码转换为结构化表示，实现从"文本变更"向"语义变更"的转化，为后续分析提供可靠基础\cite{Kimmig2018TreeSitter}。

（2）影响分析模型设计。构建基于调用关系与依赖关系的影响分析模型。调用关系用于描述函数或方法之间的调用路径；依赖关系分析涵盖模块依赖、类依赖及数据依赖等，用于识别因接口或数据变化引发的连锁影响。通过结合调用关系传播路径与依赖关系分析，构建完整的变更影响传播链路，实现从局部变更到全局影响范围的推导\cite{Horwitz1990,Nielson2010Static}。

（3）缺陷模式挖掘与融合机制。引入历史缺陷数据，对典型缺陷进行模式化建模。通过分析历史缺陷记录，总结常见缺陷特征，构建缺陷模式库。在实际应用中，将代码变更与缺陷模式进行匹配，识别潜在风险点并指导测试重点的确定。

（4）测试规划与用例生成方法。设计测试规划Agent进行统一决策，通过融合多源信息生成测试策略。结合大语言模型与Prompt工程方法，将测试点转化为结构化测试用例，包括测试步骤、输入数据及预期结果等\cite{Wang2021CodeT5}。

（5）多智能体协同框架设计。构建多智能体协同框架，各Agent分别负责代码变更分析、影响分析、测试规划等功能，通过统一状态管理机制进行调度，实现信息闭环与任务联动。

\section{本文组织架构}
本文共分为六章。第一章为绪论，介绍研究背景与意义，分析国内外研究现状，阐述本文的主要研究内容与组织架构。第二章为相关技术基础，介绍代码分析技术、影响分析方法、大语言模型与Prompt工程、多智能体系统框架等技术。第三章为系统总体设计，阐述系统的总体架构与设计原则。第四章为关键技术与实现，详细阐述各核心模块的设计与实现。第五章为实验与结果分析，通过实验验证所提方法的有效性。第六章为总结与展望，总结全文并展望未来工作。
